\RequirePackage{plautopatch} % 海外製パッケージとの互換性を保つパッチ（\documentclass より前に書く）
\documentclass[uplatex,dvipdfmx]{jlreq}

% --- 数理系必須パッケージ ---
\usepackage{amsmath,amssymb,amsthm,mathtools,bm,physics}
\usepackage{graphicx}
\usepackage[style=numeric,backend=biber]{biblatex}
\addbibresource{references.bib}
\usepackage{hyperref}
\usepackage{pxjahyper}

% --- 定理環境 ---
\theoremstyle{definition}
\newtheorem{theorem}{定理}[section]
\newtheorem{definition}[theorem]{定義}
\newtheorem{lemma}[theorem]{補題}
\renewcommand{\proofname}{\textbf{証明}}

% --- 表紙情報（ここを書き換える） ---
\title{レポートのタイトル}
\author{学籍番号：000000 \quad 氏名：学生 太郎}
\date{\today}

\begin{document}
\maketitle

% 目次（不要ならコメントアウト）
\tableofcontents
\clearpage

\section{序論}
ここに序論を記述する。研究の背景や目的など。

\section{本論}
ここに本論を記述する。

\subsection{理論}
必要な数式などを記述する。

\subsection{解析結果}
図表などをここに配置する。

\section{結論}
ここに結論を記述する。

% 謝辞（必要ならコメントアウトを外す）
%\section*{謝辞}
%指導教員や協力者への感謝を記述する。

% 参考文献
\printbibliography[title=参考文献]

\end{document}
